%---------------------------------------------------------------------------------------
% File:                 phd_thesis_template.tex
% Original Author:      Ivan Colucci
% Revised & Adapted:    Carlos Pavon
% Created:              January 2020
% Last Modified:        22/05/2025

% Description:
%   A customizable LaTeX template for PhD theses at the
%   Università degli Studi di Salerno within the
%   National PhD on Photovoltaics.

% Usage:
%   Free for non-commercial academic use. 
%   Modifiable and shareable within the National PhD on Photovoltaics program.
%   Recommended for official submissions and archiving purposes.

% Change Log:
%   [v1.0]  Initial release with institutional header.................... IC
%   [v1.1]  Adapted for National PhD program, formatting adjustments, 
%           comments, and package management............................. CP
%   [v1.2]  ............................................................. XX

% Compilation Requirements:
%   Ensure the following files and directories are present for successful compilation:
%   - Figures/: image.png           (Add images - organize in subfolders if needed)
%   - Sections/: Introduction.tex   (Add more sections as needed)
%   - main.tex                  
%   - logo_unisa.png            
%   - references.bib            

% Notes:
%   This template includes dummy text (Lorem Ipsum) and placeholders 
%   for images and tables, serving as examples for usage. 
%
%   Wishing all PhD candidates a productive and smooth writing experience!
%---------------------------------------------------------------------------------------

% DEFINES OVERALL LAYOUT AND STRUCTURE
\documentclass[a4paper, 12pt, twoside]{book}

% BASIC DOCUMENT SETUP AND ENCODING
\usepackage[utf8]{inputenc}         % Input encoding for special characters
\usepackage[T1]{fontenc}            % Hyphenation and character display
\usepackage[margin=3cm, bindingoffset=1cm]
            {geometry}              % Sets page margins
\linespread{1.5}                    % Sets line spacing

% MATHEMATICS & SYMBOLS
\usepackage{amsmath}                % Features for mathematical equations
\usepackage{amssymb}                % Additional mathematical symbols
\usepackage{amsfonts}               % Provides additional math fonts
\usepackage{gensymb}                % Symbols like degree, permille, etc.
\usepackage{latexsym}               % Provides additional LaTeX symbols

% REFERENCING AND LINKS
\usepackage{cite}                           % Citation handling and formatting
\usepackage[hidelinks]{hyperref}            % Hyperlinks within the document
    \newcommand{\figref}[1]{Fig.~\ref{#1}}  % References for Figures
    \newcommand{\tabref}[1]{Tab.~\ref{#1}}  % References for Tables
    \newcommand{\equref}[1]{Eq.~\ref{#1}}   % References for Equations

% GRAPHICS AND IMAGES
\usepackage{graphicx}               % For including images
\usepackage{svg}                    % Enables inclusion of SVG images
\usepackage{float}                  % Provides enhanced control over float placement

% CAPTIONS PERSONALIZATION
\usepackage[font=small,skip=1pt,
            labelfont=sf,textfont=normalfont]
            {caption}               % Customizes caption appearance
\usepackage{subcaption}             % create subfigures and subtables

% ALGORITHMS AND CODE LISTINGS
\usepackage{algorithmic}            % Tools for typesetting pseudocode algorithms
\usepackage{listings}               % Formatting of source code listings

% TEXT, COLORS, AND TABLES
\usepackage{textcomp}               % Additional text symbols (€, ©, ™)
\usepackage{xcolor}                 % Adds color capabilities to your document
\usepackage{csquotes}               % Provides context-sensitive quotation marks
\usepackage{booktabs}               % Enhances tables with professional-looking rules
\usepackage{multirow}               % Allows table cells to span multiple rows
\usepackage{colortbl}               % Color of individual rows, columns, or cells in tables
\usepackage{array}                  % Functionality of 'array' and 'tabular'
\usepackage{tabularray}             % Modern package for creating complex tables

% LAYOUT AND SPACING
\usepackage{indentfirst}            % Indents the first paragraph of each chapter or section
\setlength{\parindent}{1cm}         % Sets the default paragraph indentation

% HEADERS, FOOTERS, AND CHAPTER STYLE
\usepackage{fancyhdr}               % Customizes headers and footers
\setlength{\headheight}{10pt}       % Required by fancyhdr to prevent warnings
\addtolength{\topmargin}{-2.5pt}    % Adjusts top margin slightly

% UTILITIES AND MISCELLANEOUS
\usepackage{blindtext}              % Generates dummy Lorem Ipsum-like text
\usepackage{tikz}                   % Creating vector graphics
\usepackage{comment}                % Comment out large blocks of text
\usepackage{import}                 % Importing content from other files
\usepackage{xkeyval}                % Provides key-value parsing capabilities

% SPECIFIC TECHNICAL PACKAGES
\usepackage[american]{circuitikz}   % Draws electrical circuits using TikZ

% ORCID
\usepackage{orcidlink}              % Provides a command to link to an ORCID profile

% PATH FOR GRAPHICS
\graphicspath{ {./Figures/} }       % Specifies the default directory for figures

%% CUSTOM COMMANDS AND DEFINITIONS %%

% Bibtex
\def\BibTeX{{\rm B\kern-.05em{\sc i\kern-.025em b}\kern-.08em
  T\kern-.1667em\lower.7ex\hbox{E}\kern-.125emX}}   % BibTeX formatting definition

% SPICE in listings
\lstdefinelanguage{Spice}{
  morekeywords={.meas, DC, FIND, WHEN, MAX, .step, PARAM},
  sensitive=false,                                  % Keywords are not case-sensitive
  morecomment=[l]{*},                               % Single-line comments start with '*'
  morestring=[b]"                                   % Strings are enclosed in double quotes
}

% ORCID Logo Setup
\usepackage{scalerel}                               % For scaling the ORCID icon
\usetikzlibrary{svg.path}                           % Enables TikZ to draw SVG paths
\definecolor{orcidlogocol}{HTML}{A6CE39}            % Official ORCID green color
        \tikzset{
          orcidlogo/.pic={
            \fill[orcidlogocol] svg{M256,128c0,70.7-57.3,128-128,128C57.3,256,0,198.7,0,128C0,57.3,57.3,0,128,0C198.7,0,256,57.3,256,128z};
            \fill[white] svg{M86.3,186.2H70.9V79.1h15.4v48.4V186.2z}
                         svg{M108.9,79.1h41.6c39.6,0,57,28.3,57,53.6c0,27.5-21.5,53.6-56.8,53.6h-41.8V79.1z 
                             M124.3,172.4h24.5c34.9,0,42.9-26.5,42.9-39.7c0-21.5-13.7-39.7-43.7-39.7h-23.7V172.4z}
                         svg{M88.7,56.8c0,5.5-4.5,10.1-10.1,10.1
                             c-5.6,0-10.1-4.6-10.1-10.1
                             c0-5.6,4.5-10.1,10.1-10.1
                             C84.2,46.7,88.7,51.3,88.7,56.8z};
          }
        }
        \newcommand\orcidicon[1]{
          \href{https://orcid.org/#1}{
            \mbox{\scalerel*{
              \begin{tikzpicture}[yscale=-1,transform shape]
                \pic{orcidlogo};
              \end{tikzpicture}
            }{|}}
          }
        }

% Acronyms and Glossaries
\usepackage[acronym]{glossaries}                    % Manages glossaries
\newacronym{sdm}{SDM}{Single Diode Model}           % Example {key}{SHORT}{Long Form} 
    % The \gls{sdm} system...                       % In-text use for singular
    % The \glspl{sdm} system...                     % In-text use for plural
% \printglossary[type=\acronymtype, title=Acronyms] % Place to show acronyms

% Dedication environment
\newenvironment{dedication}
  {\clearpage \thispagestyle{empty} \vspace*{\stretch{1}} \itshape \raggedleft}
  {\par \vspace{\stretch{3}} \clearpage}

\begin{document}

% FRONT PAGE
\begin{titlepage}
    \begin{center}
        \LARGE{\uppercase{Università degli Studi di Salerno}}\\
        \vspace{5mm}
        \large{\textbf{NATIONAL PHD IN PHOTOVOLTAICS}}\\
        \normalsize{Curriculum: Power Electronics and Control}
    \end{center}

    \vspace{10mm}
    \begin{figure}[H]
        \centering
        \includegraphics[width=0.35\textwidth]{logo_unisa}
    \end{figure}
    
    \vspace{10mm}
    \begin{center}
        {\Huge \bfseries PhD Thesis Title}\\
        \vspace{5mm}
        \Large{\textit{A Subtitle if Needed}}\\
        \vspace{10mm}
    \end{center}
    
    \noindent
    \begin{minipage}[t]{0.47\textwidth}
        {\large{Supervisor:\\ \bf Prof. Name LastName}}\\[8mm]
        {\large{Co-Supervisor:\\ \bf Prof. Name LastName}}\\[8mm]
        {\large{PhD Programm Director:\\ \bf Prof. Giovanni Spagnuolo}}
    \end{minipage}
    \hfill
    \begin{minipage}[t]{0.4\textwidth}\raggedleft
        {\large{PhD Candidate: \\ \bf Name LastName\\ Mat. 0123456789}}
    \end{minipage}
    
    \vfill
    \centering{\large XXXVIII Cycle \\ 2022 - 2025}
\end{titlepage}

\clearpage

% Dedication
\begin{dedication}

    Dedicated to ...

\end{dedication}


% Table of Contents
    \tableofcontents
    \clearpage
    \sloppy
    \hyphenpenalty=10000
    \exhyphenpenalty=10000

% Inclusion of Chapters

\chapter{Introduction}

Lorem ipsum dolor sit amet consectetur adipiscing elit. Quisque faucibus ex sapien vitae pellentesque sem placerat. In id cursus mi pretium tellus duis convallis. Tempus leo eu aenean sed diam urna tempor. Pulvinar vivamus fringilla lacus nec metus bibendum egestas. Iaculis massa nisl malesuada lacinia integer nunc posuere. Ut hendrerit semper vel class aptent taciti sociosqu. Ad litora torquent per conubia nostra inceptos himenaeos~\cite{Topic1_Author_Year}.

\section{Example Section Title}

Lorem ipsum dolor sit amet consectetur adipiscing elit. Quisque faucibus ex sapien vitae pellentesque sem placerat. In id cursus mi pretium tellus duis convallis. Tempus leo eu aenean sed diam urna tempor. Pulvinar vivamus fringilla lacus nec metus bibendum egestas. Iaculis massa nisl malesuada lacinia integer nunc posuere \gls{sdm}. Ut hendrerit semper vel class aptent taciti sociosqu. Ad litora torquent per conubia nostra inceptos himenaeos.

\subsection{Example Subsection Title}

Donec ex turpis, eleifend at est nec, dictum egestas neque. Maecenas mollis tortor mauris, ut mattis purus volutpat sed. Vestibulum erat velit, auctor vitae felis et, tempor rhoncus sem. Etiam ut metus ante. Sed in tortor tincidunt, fermentum risus vitae, ultricies nibh. Vivamus libero ante, mollis quis leo non, pulvinar efficitur odio. Duis tincidunt dolor at nisl viverra posuere. Nullam finibus scelerisque velit, vel volutpat dolor auctor quis~\cite{Topic2_Author_Year}.

Lorem ipsum dolor sit amet consectetur adipiscing elit. Quisque faucibus ex sapien vitae pellentesque sem placerat. In id cursus mi pretium tellus duis convallis. Tempus leo eu aenean sed diam urna tempor. Pulvinar vivamus fringilla \gls{sdm} lacus nec metus bibendum egestas. Iaculis massa nisl malesuada lacinia integer nunc posuere. Ut hendrerit semper vel class aptent taciti sociosqu. Ad litora torquent per conubia nostra inceptos himenaeos~\figref{fig:placeholder_example}.

    \begin{figure}[hbtp!]
        \centering
        \includegraphics[width=0.5\linewidth]{Figures/image.png}
        \caption{Example caption for a figure.\\ \textit{Additional explanations can be added}}
        \label{fig:placeholder_example}
    \end{figure}

Phasellus posuere consequat lectus eu fringilla. In interdum fermentum quam in tempus. Cras at ligula urna. Donec non lorem elit. Etiam rhoncus nec neque quis mollis. Class aptent taciti sociosqu ad litora torquent per conubia nostra, per inceptos himenaeos. Etiam ac nisl dolor. Aenean eu justo magna. Fusce tristique nisl aliquet tempor rutrum. Curabitur nec semper sapien. Suspendisse maximus dapibus nunc ac efficitur. Aenean mollis ornare urna eu hendrerit. Suspendisse eu lorem egestas, laoreet orci eget, finibus ligula. Vestibulum vel dictum diam.

Donec ex turpis, eleifend at est nec, dictum egestas neque. Maecenas mollis tortor mauris, ut mattis purus volutpat sed. Vestibulum erat velit, auctor vitae felis et, tempor rhoncus sem. Etiam ut metus ante. Sed in tortor tincidunt, fermentum risus vitae, ultricies nibh. Vivamus libero ante, mollis quis leo non, pulvinar efficitur odio. Duis tincidunt dolor at nisl viverra posuere. Nullam finibus scelerisque velit, vel volutpat dolor auctor quis.

Lorem ipsum dolor sit amet consectetur adipiscing elit. Quisque faucibus ex sapien vitae pellentesque sem placerat. In id cursus mi pretium tellus duis convallis. Tempus leo eu aenean sed diam urna tempor. Pulvinar vivamus fringilla lacus nec metus bibendum egestas. Iaculis massa nisl malesuada lacinia integer nunc posuere. Ut hendrerit semper vel class aptent taciti sociosqu. Ad litora torquent per \gls{sdm} conubia nostra inceptos himenaeos~\tabref{tab:placeholder_table}.

    \begin{table}[ht]
    \centering
        \begin{tabular}{|l|l|l|l|l|}
            \hline
            \textbf{Parameter} & \textbf{Unit} & \textbf{Test 1} & \textbf{Test 2} & \textbf{Test 3} \\ \hline
            Input Voltage      & V             & 12.0            & 11.8            & 12.2            \\ \hline
            Output Current     & A             & 1.5             & 1.6             & 1.4             \\ \hline
            Efficiency         & \%            & 92.3            & 91.8            & 93.1            \\ \hline
            Temperature Rise   & °C            & 32              & 34              & 31              \\ \hline
        \end{tabular}
        \caption{Performance comparison of the table caption.}
        \label{tab:placeholder_table}
    \end{table}

Aliquam dui justo, eleifend non tempus quis, fringilla nec ante. Fusce varius nulla sit amet erat dapibus gravida. Nam vitae turpis nunc. Nulla convallis ultricies dolor, id luctus mi rhoncus vel. Vivamus pharetra laoreet massa eget aliquet. Donec in diam in ligula accumsan aliquet. Morbi fringilla non massa eget pulvinar. Aenean sit amet posuere felis, sed iaculis tortor. Sed cursus, nibh sed ullamcorper finibus, purus libero malesuada leo, et porta purus orci a nisi. Nam sit amet justo tortor. Donec ac ex aliquam orci facilisis porttitor. Morbi finibus, nibh at mattis sollicitudin, nunc nunc lobortis justo, sed vehicula neque urna a odio \equref{eq:placeholder_equation}.

    \begin{equation}
        i_{pv}(t) = i_{ph}(t) - i_{D}(t) - i_{R_{sh}}(t)
        \label{eq:placeholder_equation}
    \end{equation}

Nam sed odio nulla. Cras et velit in ex tempus cursus. Etiam scelerisque sapien at pulvinar placerat. Cras consequat nunc ligula, sit amet facilisis mauris fringilla vel. Sed blandit porttitor elit non laoreet. Duis placerat purus eget magna varius suscipit. Aliquam est tortor, tincidunt vitae odio ut, egestas convallis sapien. Donec ex turpis, eleifend at est nec, dictum egestas neque. Maecenas mollis tortor mauris, ut mattis purus volutpat sed. Vestibulum erat velit, auctor vitae felis et, tempor rhoncus sem. Etiam ut metus ante. Sed in tortor tincidunt, fermentum risus vitae, ultricies nibh. Vivamus libero ante, mollis quis leo non, pulvinar efficitur odio. Duis tincidunt dolor at nisl viverra posuere. Nullam finibus scelerisque velit, vel volutpat dolor auctor quis. % Recommended practice for maintaining a well-organized main.tex file

\chapter{Literature Review}
\blindtext[7] % Generates Lorem Ipsum

\chapter{Methodology}
\blindtext[5] % Generates Lorem Ipsum

\chapter{Results and Discussion}
\blindtext[7] % Generates Lorem Ipsum

\chapter{Conclusion}
\blindtext[3] % Generates Lorem Ipsum

% Acknowledgements
\chapter*{Acknowledgements}
\thispagestyle{empty}
\markboth{ACKNOWLEDGEMENTS}{}
\addcontentsline{toc}{chapter}{Acknowledgements}
\blindtext[1] % Generates Lorem Ipsum

% Bibliography
\clearpage
\bibliographystyle{ieeetr}
\bibliography{references.bib} % Make sure this file exists

% Lists of Figures and Tables
\clearpage
\listoffigures
\addcontentsline{toc}{chapter}{List of Figures}

\clearpage
\listoftables
\addcontentsline{toc}{chapter}{List of Tables}

\clearpage
\end{document}